% from ega4-12-fr.tex
\oldpage[IV-3]{175}
qu'on est ramené au cas où $Y=\operatorname{Spec}(A)$ est le spectre d'un
anneau de valuation discrète, $y$ étant le point fermé et $y'$ le point
générique de $Y$. Si $t$ est une uniformisante de $A$, l'hypothèse que
$\mathcal F$ est $f$-plat entraîne que $t$ est $\mathcal F$-régulier
($\mathbf 0_{\mathrm I}$, 6.3.4). Par hypothèse aucun des cycles premiers
associés immergés de $\mathcal F/t\mathcal F$ ne contient $x$. Alors, il
résulte de (3.4.4) que $x$ n'appartient à aucun cycle premier immergé associé à
$\mathcal F$, et il en est donc de même de toute générisation de $x$. En
particulier, $x'$ n'appartient à aucun cycle premier immergé associé à
$\mathcal F$, ni \emph{a fortiori} à aucun de ceux associés à
$\mathcal F_{y'}$ ($X_{y'}$ étant un sous-préschéma induit sur un ouvert de
$X$).
\end{env}

\begin{env}[12.1.1.2]
\label{IV.12.1.1.2-fr}
Envisageons ensuite les cas (iv) et (v) ((vi) n'est qu'un cas particulier de
(v)). On procède comme ci-dessus (en utilisant (6.7.1) au lieu de (4.2.7)) et
l'on est ramené au cas où $Y$ est le spectre d'un anneau de valuation discrète
$A$.

L'anneau $C=\mathcal O_{X,x}$ est alors localisé d'une $A$-algèbre de type
fini, donc est caténaire (5.6.4), et par hypothèse le $C$-module
$M_x/tM_x$ vérifie $(S_k)$ et est équidimensionnel (resp. on a
$\operatorname{coprof}(M_x/tM_x)\leq k$)~; on déduit donc de (5.12.2) (resp.
de ($\mathbf 0$, 16.4.10)) que $M_x$ vérifie $(S_k)$ et est équidimensionnel
(resp. que $\operatorname{coprof}(M_x)\leq k$) puisque $t$ est $M_x$-régulier
et appartient à l'idéal maximal de $C$. D'où les conclusions puisque $x'$ est
une générisation de $x$ et que $(\mathcal F_{y'})_{x'}=\mathcal F_{x'}$.
\end{env}

\begin{env}[12.1.1.3]
\label{IV.12.1.1.3-fr}
Passons aux cas (vii) et (viii). On peut évidemment remplacer $Y$ par le
sous-schéma intègre ayant $\overline{\{y'\}}$ pour espace sous-jacent, et $X$
par $f^{-1}(\overline{\{y'\}})$, sans changer les fibres en $y$ et $y'$, donc
on peut supposer $A$ intègre, de corps des fractions $K=k(y')$. On sait
((4.5.11) et (4.7.8)) qu'il existe une extension finie $K'$ de $K$ telle que,
pour le $(\mathcal O_X\otimes_AK')$-Module
$\mathcal F'=\mathcal F\otimes_AK'$, les cycles premiers associés soient
géométriquement irréductibles et que les longueurs géométriques de
$\mathcal F'$ aux points maximaux $z'_i$ de son support soient respectivement
égales aux longueurs de $\mathcal F'_{z'_i}$ en ces points~; il revient donc
au même de dire que $\mathcal F$ est géométriquement réduit (resp.
géométriquement ponctuellement intègre) en un point $x'$ de
$X_{y'}=X\otimes_AK$, ou de dire que $\mathcal F'$ est réduit (resp. intègre)
aux points de $X\otimes_AK'$ au-dessus de $x'$, compte tenu de (4.2.7). Soit
$A'$ une $A$-algèbre finie dont $K'$ est le corps des fractions~; il résulte
de (4.7.11) qu'en tout point de $X\otimes_AA'$ au-dessus de $x$,
$\mathcal F\otimes_AA'$ a la propriété (vii) (resp. (viii)). Remplaçant $A$
par $A'$ et $X$ par $X\otimes_AA'$, on voit qu'on peut se borner à prouver que
$\mathcal F$ est réduit (resp. intègre) en toute générisation $x'$ de $x$. On
procède alors comme dans (12.1.1.1), en utilisant cette fois (4.7.11), et l'on
est encore ramené au cas où $A$ est un anneau de valuation discrète, $y$ étant
le point fermé, $y'$ le point générique de $Y$, et l'uniformisante $t$ de $A$
étant un élément $M$-régulier. Comme par hypothèse
$\mathcal F/t\mathcal F$ est réduit (resp. intègre) au point $x$, il résulte de
(3.4.6) (resp. (3.4.5)) que $\mathcal F$ est réduit (resp. intègre) au point
$x$, donc aussi dans un voisinage de $x$, et en particulier au point $x'$, ce
qui achève la démonstration dans les cas (vii) et (viii).
\end{env}

\begin{env}[12.1.1.4]
\label{IV.12.1.1.4-fr}
Reste à examiner les cas (i) et (ii). Remplaçant $Y$ par le sous-schéma intègre
ayant $\overline{\{y'\}}$ pour espace sous-jacent, on peut se borner au cas où
$Y$ est irréductible et $y'$ son point générique. Soit $z'$ un point générique
d'un cycle premier associé à $\mathcal F_{y'}$
